\documentclass{ximera}

\title{Vraagstukken potentiële energie, potentiaal en spanning: spanning en arbeid}
\author{Daan Lipkens}

\begin{document}
\begin{abstract}
    Oefeningen op het verband tussen spanning en arbeid ($W = Q \cdot U$), van makkelijk naar moeilijk.
\end{abstract}
\maketitle

\section{Oefeningen spanning en arbeid}

%———————————————————————————————————————
% OEF 1: W = Q*U
\begin{exercise}
    Een lading $Q = 2{,}0 \cdot 10^{-6} \, \mathrm{C}$ doorloopt een spanningsverschil van $500 \, \mathrm{V}$. Bereken de arbeid die het elektrisch veld op de lading verricht.

    \[ W = \answer[tolerance=0.01,onlineshowanswerbutton]{1.0} \cdot 10^{-3} \, \mathrm{J} \]

    \begin{solution}\nl
        \underline{Gegevens:}
        \begin{align*}
            & Q = 2{,}0 \cdot 10^{-6} \, \mathrm{C} & & U = 500 \, \mathrm{V}
        \end{align*}

        \underline{Gevraagd:} $W$

        \underline{Oplossing:}
        \begin{align*}
            W = Q \cdot U = 2{,}0 \cdot 10^{-6} \, \mathrm{C} \cdot 500 \, \mathrm{V} = 1{,}0 \cdot 10^{-3} \, \mathrm{J}
        \end{align*}
    \end{solution}
\end{exercise}
%———————————————————————————————————————

%———————————————————————————————————————
% OEF 2: U = W/Q
\begin{exercise}
    Het elektrisch veld verricht een arbeid van $2{,}0 \cdot 10^{-3} \, \mathrm{J}$ op een lading $Q = +5{,}0 \cdot 10^{-6} \, \mathrm{C}$ die van punt $1$ naar punt $2$ beweegt. Bereken de spanning $U_{12}$ tussen die twee punten.

    \[ U_{12} = \answer[tolerance=1,onlineshowanswerbutton]{400} \, \mathrm{V} \]

    \begin{solution}\nl
        \underline{Gegevens:}
        \begin{align*}
            & W = 2{,}0 \cdot 10^{-3} \, \mathrm{J} & & Q = +5{,}0 \cdot 10^{-6} \, \mathrm{C}
        \end{align*}

        \underline{Gevraagd:} $U_{12}$

        \underline{Oplossing:}
        \begin{align*}
            U_{12} = \frac{W}{Q} = \frac{2{,}0 \cdot 10^{-3} \, \mathrm{J}}{5{,}0 \cdot 10^{-6} \, \mathrm{C}} = 4{,}0 \cdot 10^{2} \, \mathrm{V}
        \end{align*}
    \end{solution}
\end{exercise}
%———————————————————————————————————————

%———————————————————————————————————————
% OEF 3: Q = W/U (autobatterij)
\begin{exercise}
    Een autobatterij van $12 \, \mathrm{V}$ levert een energie van $3{,}6 \cdot 10^{3} \, \mathrm{J}$. Hoeveel lading is daarbij van de ene pool naar de andere verplaatst?

    \[ Q = \answer[tolerance=1,onlineshowanswerbutton]{300} \, \mathrm{C} \]

    \begin{solution}\nl
        \underline{Gegevens:}
        \begin{align*}
            & U = 12 \, \mathrm{V} & & W = 3{,}6 \cdot 10^{3} \, \mathrm{J}
        \end{align*}

        \underline{Gevraagd:} $Q$

        \underline{Oplossing:}
        \begin{align*}
            Q = \frac{W}{U} = \frac{3{,}6 \cdot 10^{3} \, \mathrm{J}}{12 \, \mathrm{V}} = 3{,}0 \cdot 10^{2} \, \mathrm{C}
        \end{align*}
    \end{solution}
\end{exercise}
%———————————————————————————————————————

%———————————————————————————————————————
% OEF 4 (moeilijkst): negatieve lading, teken van W
\begin{exercise}
    Een elektron beweegt van punt $1$ naar punt $2$ in een elektrisch veld. De spanning tussen die punten is $U_{12} = +80 \, \mathrm{V}$. Bereken de verandering van de kinetische energie van het elektron.

    \[ \Delta E_{\text{kin}} = \answer[tolerance=0.02,onlineshowanswerbutton]{-1.28} \cdot 10^{-17} \, \mathrm{J} \]

    Het elektron wordt: \wordChoice{
        \choice{versneld}
        \choice[correct]{afgeremd}
    }

    \begin{hint}
        Houd rekening met het teken van de lading van een elektron.
    \end{hint}

    \begin{solution}\nl
        \underline{Gegevens:}
        \begin{align*}
            & Q = -e = -1{,}60 \cdot 10^{-19} \, \mathrm{C} & & U_{12} = +80 \, \mathrm{V}
        \end{align*}

        \underline{Gevraagd:} $\Delta E_{\text{kin}}$

        \underline{Oplossing:}
        \begin{align*}
            \Delta E_{\text{kin}} = W = Q \cdot U_{12} = (-1{,}60 \cdot 10^{-19} \, \mathrm{C}) \cdot 80 \, \mathrm{V} = -1{,}28 \cdot 10^{-17} \, \mathrm{J}
        \end{align*}
        De verandering van de kinetische energie is negatief: het elektron wordt afgeremd. Dit is logisch: het elektron beweegt naar een punt met een lagere potentiaal, terwijl negatieve ladingen net naar een hogere potentiaal worden aangetrokken.
    \end{solution}
\end{exercise}
%———————————————————————————————————————

\end{document}